\documentclass[12pt, a4paper]{article}

% ── Encoding & Language ──────────────────────────────────────────────────────
\usepackage[utf8]{inputenc}
\usepackage[T1]{fontenc}
\usepackage[italian]{babel}

% ── Layout ───────────────────────────────────────────────────────────────────
\usepackage[top=2.5cm, bottom=2.5cm, left=3cm, right=3cm]{geometry}
\usepackage{setspace}
\onehalfspacing

% ── Typography ────────────────────────────────────────────────────────────────
\usepackage{lmodern}
\usepackage{microtype}

% ── Tables ───────────────────────────────────────────────────────────────────
\usepackage{booktabs}
\usepackage{longtable}
\usepackage{array}
\usepackage{tabularx}
\usepackage{multirow}
\usepackage{xcolor}
\definecolor{headerblue}{RGB}{41, 82, 148}
\definecolor{rowgray}{RGB}{240, 240, 240}

% ── Graphics & Diagrams ───────────────────────────────────────────────────────
\usepackage{graphicx}
\usepackage{tikz}

% ── Hyperlinks ────────────────────────────────────────────────────────────────
\usepackage[colorlinks=true, linkcolor=headerblue, urlcolor=headerblue, citecolor=headerblue]{hyperref}

% ── Misc ─────────────────────────────────────────────────────────────────────
\usepackage{enumitem}
\usepackage{fancyhdr}
\usepackage{lastpage}
\usepackage{titlesec}
\usepackage{caption}

% ── Header / Footer ──────────────────────────────────────────────────────────
\pagestyle{fancy}
\fancyhf{}
\fancyhead[L]{\small Solco Fiscale – Software Requirements Specification}
\fancyhead[R]{\small v0.1 – \today}
\fancyfoot[C]{\small Pagina \thepage\ di \pageref{LastPage}}
\renewcommand{\headrulewidth}{0.4pt}

% ── Section Formatting ───────────────────────────────────────────────────────
\titleformat{\section}{\large\bfseries\color{headerblue}}{{\thesection}}{1em}{}[\titlerule]
\titleformat{\subsection}{\normalsize\bfseries}{\thesubsection}{1em}{}
\titleformat{\subsubsection}{\normalsize\itshape}{\thesubsubsection}{1em}{}

% ── Custom Commands ───────────────────────────────────────────────────────────
% Requirement entry: \req{ID}{Descrizione}
\newcommand{\req}[2]{%
  \noindent\textbf{[#1]} \quad #2
}

\begin{document}

% ── Title Page ────────────────────────────────────────────────────────────────
\begin{titlepage}
  \centering
  \vspace*{2cm}
  {\Huge\bfseries\color{headerblue} Solco Fiscale\par}
  \vspace{0.5cm}
  {\large\itshape Sistema di Gestione del Catasto e Calcolo delle Imposte\par}
  \vspace{2cm}
  {\LARGE Software Requirements Specification\par}
  \vspace{0.5cm}
  {\large Versione 0.1 -- Bozza\par}
  \vspace{3cm}
  \begin{tabular}{ll}
    \textbf{Cliente:}   & Agenzia Regionale, Emilia-Romagna \\[4pt]
    \textbf{Progetto:}  & Solco Fiscale                     \\[4pt]
    \textbf{Data:}      & \today                            \\[4pt]
    \textbf{Stato:}     & Bozza iniziale                    \\[4pt]
    \textbf{Standard:}  & IEEE 830-1998                     \\
  \end{tabular}
  \vfill
  {\small Documento riservato -- uso interno al gruppo di sviluppo}
\end{titlepage}

% ── Change History ────────────────────────────────────────────────────────────
\section*{Storico delle Revisioni}
\begin{tabularx}{\textwidth}{lllX}
  \toprule
  \textbf{Versione} & \textbf{Data} & \textbf{Autore} & \textbf{Descrizione} \\
  \midrule
  0.1 & \today & Gruppo & Prima bozza -- struttura e requisiti preliminari \\
  \bottomrule
\end{tabularx}
\vspace{1cm}

% ── Table of Contents ─────────────────────────────────────────────────────────
\tableofcontents
\newpage

% =============================================================================
\section{Introduzione}
% =============================================================================

\subsection{Scopo del documento}

Questo documento costituisce la \textit{Software Requirements Specification} (SRS) del sistema
\textbf{Solco Fiscale}, redatta secondo lo standard IEEE~830-1998.

Lo scopo è definire in modo chiaro e non ambiguo i requisiti funzionali e non funzionali
concordati tra il committente (Agenzia Regionale Emilia-Romagna) e il gruppo di sviluppo,
costituendo il contratto tecnico di riferimento per tutto il ciclo di vita del progetto.

Il documento è destinato a:
\begin{itemize}
  \item Il committente, per approvazione e validazione dei requisiti.
  \item Il gruppo di sviluppo, come riferimento per la progettazione e l'implementazione.
  \item Il responsabile del collaudo, per la definizione dei criteri di accettazione.
\end{itemize}

\subsection{Ambito del sistema (\textit{Scope})}

\textbf{Solco Fiscale} è un sistema web per la gestione delle imposte sui terreni agricoli
catastali nella regione Emilia-Romagna. Il sistema sostituirà l'applicazione legacy
attualmente in uso (sviluppata in Visual Basic, eseguita su macchine Windows~7 in rete locale).

\medskip\noindent\textbf{Obiettivi principali:}
\begin{enumerate}
  \item Automatizzare l'importazione dei dati catastali forniti mensilmente dal Catasto.
  \item Calcolare le imposte fondiarie utilizzando l'algoritmo esistente (da integrare).
  \item Gestire il flusso di approvazione mensile e l'invio delle notifiche ai contribuenti.
  \item Mantenere uno storico completo delle operazioni e garantire la tracciabilità.
\end{enumerate}

\medskip\noindent\textbf{Fuori ambito (versione 1.0):}
\begin{itemize}
  \item Portale pubblico per i contribuenti con autenticazione SPID.
  \item Applicazione mobile.
  \item Integrazione diretta in tempo reale con le banche dati del Catasto nazionale.
  \item Analisi avanzate o reportistica BI.
\end{itemize}

\subsection{Definizioni, acronimi e abbreviazioni}

\begin{longtable}{lp{10cm}}
  \toprule
  \textbf{Termine} & \textbf{Definizione} \\
  \midrule
  \endhead
  SRS       & Software Requirements Specification \\
  Catasto   & Ufficio del Catasto; ente che fornisce i dati fondiari mensili \\
  Dato catastale & Record che descrive un terreno: codice fiscale del proprietario, superficie, tipo di terreno, Comune, Provincia \\
  Imposta fondiaria & Tributo calcolato sul valore catastale del terreno \\
  Algoritmo legacy & Algoritmo di calcolo delle imposte attualmente implementato in Visual Basic \\
  Operatore & Utente con privilegi limitati (inserimento dati, visualizzazione) \\
  Responsabile di Area & Utente con autorità di approvazione delle operazioni critiche \\
  Amministratore & Utente IT con accesso completo al sistema \\
  PDF di notifica & Documento generato dal sistema e inviato ai contribuenti \\
  Flusso mensile & Sequenza: importazione dati $\to$ calcolo imposte $\to$ approvazione $\to$ invio notifiche \\
  SPID      & Sistema Pubblico di Identità Digitale (identità digitale italiana) \\
  RBAC      & Role-Based Access Control \\
  \bottomrule
\end{longtable}

\subsection{Riferimenti}

\begin{enumerate}
  \item IEEE Std 830-1998 -- \textit{IEEE Recommended Practice for Software Requirements Specifications}
  \item Verbale del colloquio con il cliente -- 30 marzo 2026 (\texttt{dondi\_requirment.md})
  \item Note di analisi del business analyst (\texttt{giansoldati\_requirment.md})
  \item Sintesi dei requisiti di sistema (\texttt{menozzi\_requirment.md})
  \item Template Excel del Catasto (da acquisire)
  \item Documentazione dell'algoritmo legacy Visual Basic (da acquisire)
\end{enumerate}

\subsection{Panoramica del documento}

La Sezione~2 descrive il sistema a livello generale: contesto, utenti, vincoli e dipendenze.
La Sezione~3 dettaglia i requisiti funzionali e non funzionali.
L'Appendice riporta i modelli \textit{as-is} e \textit{to-be} del flusso operativo.

% =============================================================================
\section{Descrizione Generale}
% =============================================================================

\subsection{Prospettiva del prodotto}

Solco Fiscale si inserisce nel contesto di un'agenzia regionale che gestisce l'imposizione
fiscale sui terreni agricoli. Il sistema sostituisce un'applicazione desktop legacy e introduce:
\begin{itemize}
  \item Un'interfaccia web accessibile anche in modalità \textit{smart working}.
  \item Un flusso di importazione dati semiautomatico in sostituzione dell'inserimento manuale.
  \item Un motore di calcolo integrato con l'algoritmo legacy esistente.
\end{itemize}

\subsubsection{Interfaccia di sistema}

Il sistema sarà erogato come applicazione web (client-server). La scelta dell'infrastruttura
(Azure cloud vs.\ server on-premise) è ancora aperta e verrà concordata con il committente;
il sistema deve essere progettato per supportare entrambe le opzioni senza modifiche
architetturali sostanziali.

\subsubsection{Interfaccia utente}

Interfaccia web accessibile da browser moderni (Chrome, Edge, Firefox). Non è richiesta
l'installazione di software sul client. L'interfaccia deve essere sufficientemente simile
all'applicazione legacy da minimizzare il costo di riaddestramento degli operatori.

\subsubsection{Interfaccia hardware}

Gli utenti utilizzano attualmente PC con Windows~7. La soluzione web elimina la dipendenza
dal sistema operativo client. Almeno un PC è dotato di connessione a Internet (necessaria
per l'invio delle email).

\subsubsection{Interfaccia software}

\begin{itemize}
  \item \textbf{Algoritmo legacy:} Da integrare come componente o servizio separato;
        l'interfaccia esatta (chiamata diretta, microservizio, riscrittura) è da definire.
  \item \textbf{Importazione Excel:} Il sistema deve leggere i file Excel forniti dal Catasto
        nel formato attuale (schema da acquisire).
  \item \textbf{Servizio email:} Invio tramite SMTP (server da concordare con il committente).
  \item \textbf{Generazione PDF:} Libreria lato server per la produzione dei PDF di notifica.
\end{itemize}

\subsubsection{Interfaccia di comunicazione}

\begin{itemize}
  \item HTTPS per tutta la comunicazione web.
  \item SMTP (con TLS) per l'invio delle email ai contribuenti.
  \item Eventuale VPN per accesso remoto all'infrastruttura on-premise.
\end{itemize}

\subsubsection{Requisiti di memoria e archiviazione}

Il volume dei dati è contenuto: circa 20--30 record catastali al mese con storico annuale.
La stima del database è nell'ordine di poche decine di GB nel lungo periodo.
Backup giornalieri richiesti (modalità da definire in base alla scelta infrastrutturale).

\subsubsection{Inizializzazione, backup e ripristino}

\begin{itemize}
  \item Backup giornaliero automatico del database.
  \item Procedura di ripristino documentata e testata prima del go-live.
  \item Migrazione dei dati storici dal sistema legacy (piano da definire).
\end{itemize}

\subsection{Funzioni del prodotto (sintesi)}

\begin{enumerate}
  \item \textbf{Autenticazione e gestione degli accessi} -- login con ruolo assegnato.
  \item \textbf{Importazione dati catastali} -- da file Excel con validazione e workflow di approvazione.
  \item \textbf{Inserimento manuale} -- fallback per dati non importabili automaticamente.
  \item \textbf{Calcolo delle imposte} -- tramite algoritmo legacy integrato, eseguito mensilmente.
  \item \textbf{Approvazione e invio notifiche} -- flusso: calcolo $\to$ revisione manager $\to$ generazione PDF $\to$ invio email.
  \item \textbf{Gestione dello storico} -- tracciamento di tutte le modifiche ai dati catastali.
  \item \textbf{Audit log} -- registro delle operazioni utente, visibile solo agli amministratori.
  \item \textbf{Ricerca e filtri} -- per territorio, codice fiscale, Provincia.
\end{enumerate}

\subsection{Caratteristiche degli utenti}

\begin{tabularx}{\textwidth}{lXX}
  \toprule
  \textbf{Ruolo} & \textbf{Profilo} & \textbf{Competenze tecniche} \\
  \midrule
  Operatore & 4--5 impiegati, uso quotidiano del sistema per inserimento e verifica dati & Basse; formati Office, web browser \\
  \addlinespace
  Responsabile di Area & 1--2 responsabili, usano il sistema per approvazioni mensili & Medie; familiarità con Excel e email \\
  \addlinespace
  Amministratore IT & 1 persona, gestione tecnica del sistema & Elevate; gestione server/cloud \\
  \bottomrule
\end{tabularx}

\subsection{Vincoli generali}

\begin{itemize}
  \item \textbf{Budget:} €\,50.000--60.000 (finanziamento regionale).
  \item \textbf{Algoritmo di calcolo:} L'algoritmo legacy deve essere riutilizzato senza
        alterarne la logica, salvo evidenze di malfunzionamento.
  \item \textbf{Utenti concorrenti:} Massimo 5--6 operatori + 2--3 responsabili
        simultaneamente; carico molto limitato.
  \item \textbf{Volume dati:} Circa 20 file Excel al mese, ciascuno con circa 20 righe.
  \item \textbf{Supporto post-lancio:} Da 6 a 12 mesi inclusi nel contratto.
  \item \textbf{Conformità:} I dati fiscali e i dati personali dei contribuenti devono essere
        trattati in conformità al GDPR.
\end{itemize}

\subsection{Assunzioni e dipendenze}

\begin{itemize}
  \item Il formato dei file Excel forniti dal Catasto rimane stabile nel tempo.
  \item Il cliente fornirà la documentazione (o il sorgente) dell'algoritmo legacy VB
        prima dell'inizio della fase di sviluppo.
  \item È disponibile un server SMTP per l'invio delle email.
  \item La scelta infrastrutturale (Azure vs.\ on-premise) sarà definita entro la fine
        della fase di progettazione.
  \item \textbf{[DA DECIDERE]} Strategia di migrazione dei dati storici dal sistema legacy.
\end{itemize}

\subsection{Requisiti rimandati}

I seguenti elementi sono fuori ambito nella versione 1.0 ma dovranno essere
considerati nell'architettura per non precluderne l'aggiunta futura:

\begin{itemize}
  \item Portale pubblico con autenticazione SPID per i contribuenti.
  \item Applicazione mobile per gli operatori in smart working.
  \item Integrazione con SMS per le notifiche ai contribuenti.
  \item Reportistica avanzata e cruscotti analitici.
\end{itemize}

% =============================================================================
\section{Requisiti Dettagliati}
% =============================================================================

\subsection{Interfaccia di sistema e frontend}

Le pagine/schermate principali dell'applicazione web sono:

\begin{itemize}
  \item \textbf{Login} -- autenticazione username/password con selezione ruolo.
  \item \textbf{Dashboard} -- riepilogo attività in corso (importazioni pendenti, approvazioni attese).
  \item \textbf{Gestione dati catastali} -- lista, ricerca, inserimento/modifica record.
  \item \textbf{Importazione Excel} -- upload file, anteprima dati, validazione, invio ad approvazione.
  \item \textbf{Calcolo imposte} -- avvio manuale o schedulato, stato elaborazione, riepilogo risultati.
  \item \textbf{Approvazione e invio notifiche} -- dashboard del Responsabile con lista PDF da approvare.
  \item \textbf{Storico operazioni} -- log ricercabile (solo Amministratore).
  \item \textbf{Gestione utenti} -- creazione/modifica/disattivazione account (solo Amministratore).
\end{itemize}

% ─────────────────────────────────────────────────────────────────────────────
\subsection{Requisiti Funzionali}
% ─────────────────────────────────────────────────────────────────────────────

Per ogni requisito sono indicati: identificatore, descrizione, input, output, comportamento
in caso di errore e ruoli abilitati.

% ── AUTH ─────────────────────────────────────────────────────────────────────
\subsubsection{Autenticazione e controllo accessi}

\begin{longtable}{@{}p{2.5cm}p{9.5cm}@{}}
  \toprule
  \textbf{ID} & \textbf{Descrizione} \\
  \midrule
  \endhead

  AUTH-01 &
  Il sistema deve permettere il login tramite username e password. In caso di credenziali
  errate, deve mostrare un messaggio generico senza rivelare quale campo è sbagliato.
  Dopo 5 tentativi falliti consecutivi, l'account viene temporaneamente bloccato. \\[6pt]

  AUTH-02 &
  Il sistema deve implementare RBAC con tre ruoli: \textbf{Amministratore},
  \textbf{Responsabile di Area}, \textbf{Operatore}. Ogni pagina e funzione è accessibile
  solo dai ruoli autorizzati; accessi non consentiti restituiscono HTTP 403. \\[6pt]

  AUTH-03 &
  La sessione utente deve scadere dopo 30 minuti di inattività; l'utente viene reindirizzato
  alla pagina di login senza perdita di dati non salvati (warning preventivo). \\[6pt]

  AUTH-04 &
  Solo l'Amministratore può creare, modificare, disattivare account utente. La disattivazione
  non cancella i dati storici associati all'account. \\

  \bottomrule
\end{longtable}

% ── DATA ─────────────────────────────────────────────────────────────────────
\subsubsection{Gestione dati catastali}

\begin{longtable}{@{}p{2.5cm}p{9.5cm}@{}}
  \toprule
  \textbf{ID} & \textbf{Descrizione} \\
  \midrule
  \endhead

  DATA-01 &
  Il sistema deve permettere l'importazione di file Excel nel formato fornito dal Catasto.
  Il file viene caricato dall'Operatore; il sistema esegue la validazione del formato e dei
  dati (campi obbligatori, tipi, range). L'esito della validazione (lista errori per riga)
  è mostrato prima di procedere. \\[6pt]

  DATA-02 &
  I record validati vengono messi in uno stato \textit{«in attesa di approvazione»};
  il Responsabile di Area li revisiona e può approvarli o rifiutarli (con motivazione).
  Solo i record approvati vengono confermati nel database. \\[6pt]

  DATA-03 &
  Il sistema deve permettere l'inserimento manuale di un singolo record catastale da parte
  dell'Operatore (campi: codice fiscale proprietario, Comune, Provincia, superficie in m²,
  tipo terreno). I campi obbligatori sono tutti. \\[6pt]

  DATA-04 &
  Nessun record catastale può essere eliminato fisicamente da Operatori o Responsabili.
  Solo l'Amministratore può eseguire eliminazioni fisiche. Gli Operatori possono
  \textit{disattivare} un record (eliminazione logica con tracciamento). \\[6pt]

  DATA-05 &
  Ogni modifica a un record catastale deve essere tracciata nello storico: utente,
  timestamp, campo modificato, valore precedente, valore nuovo. \\[6pt]

  DATA-06 &
  Il sistema deve permettere di allegare file (PDF, ZIP) a un record catastale.
  Dimensione massima del singolo allegato: \textbf{[DA DEFINIRE]} MB. \\[6pt]

  DATA-07 &
  Il sistema deve offrire funzioni di ricerca e filtraggio dei record per:
  codice fiscale del proprietario, Comune, Provincia, tipo terreno, stato (attivo/disattivo). \\

  \bottomrule
\end{longtable}

% ── CALC ─────────────────────────────────────────────────────────────────────
\subsubsection{Calcolo delle imposte}

\begin{longtable}{@{}p{2.5cm}p{9.5cm}@{}}
  \toprule
  \textbf{ID} & \textbf{Descrizione} \\
  \midrule
  \endhead

  CALC-01 &
  Il sistema deve integrare l'algoritmo di calcolo delle imposte fondiarie attualmente
  implementato in Visual Basic. La logica di calcolo non deve essere alterata. \\[6pt]

  CALC-02 &
  Il calcolo deve poter essere avviato manualmente dal Responsabile di Area o schedulato
  automaticamente (frequenza: mensile, giorno e ora da configurare). \\[6pt]

  CALC-03 &
  L'algoritmo opera su tutti i record catastali attivi. Il risultato per ciascun record
  comprende: importo imposta, tipo imposta, scadenza di pagamento, modalità di pagamento.
  I risultati sono salvati con riferimento al mese di calcolo. \\[6pt]

  CALC-04 &
  In caso di errore durante il calcolo (es.\ dati mancanti, errore algoritmo), il sistema
  deve registrare l'errore nell'audit log, notificare l'Amministratore e non sovrascrivere
  i risultati precedenti. \\

  \bottomrule
\end{longtable}

% ── NOTIF ─────────────────────────────────────────────────────────────────────
\subsubsection{Generazione e invio notifiche}

\begin{longtable}{@{}p{2.5cm}p{9.5cm}@{}}
  \toprule
  \textbf{ID} & \textbf{Descrizione} \\
  \midrule
  \endhead

  NOTIF-01 &
  Dopo il calcolo mensile, il sistema deve generare automaticamente un PDF di notifica per
  ciascun contribuente, contenente: dettaglio terreni, importo imposta, scadenze. \\[6pt]

  NOTIF-02 &
  I PDF generati devono essere presentati al Responsabile di Area in una dashboard di
  revisione. Il Responsabile può approvarli (singolarmente o in blocco) o rifiutarli. \\[6pt]

  NOTIF-03 &
  Dopo l'approvazione, il sistema invia i PDF via email ai rispettivi contribuenti.
  L'invio può essere avviato manualmente (pulsante «Invia») o in modo schedulato. \\[6pt]

  NOTIF-04 &
  Il sistema deve registrare per ogni email: destinatario, timestamp invio, esito
  (successo/fallimento). In caso di fallimento, l'Operatore viene notificato. \\[6pt]

  NOTIF-05 &
  I PDF approvati e inviati devono essere archiviati nel sistema e associabili al relativo
  mese di competenza per futura consultazione. \\

  \bottomrule
\end{longtable}

% ── AUDIT ─────────────────────────────────────────────────────────────────────
\subsubsection{Audit log}

\begin{longtable}{@{}p{2.5cm}p{9.5cm}@{}}
  \toprule
  \textbf{ID} & \textbf{Descrizione} \\
  \midrule
  \endhead

  AUDIT-01 &
  Il sistema deve registrare in un log tutte le operazioni significative: login/logout,
  import file, modifiche a record catastali, approvazioni, invii email, modifiche utenti.
  Ogni voce del log contiene: timestamp, utente, tipo operazione, dettaglio. \\[6pt]

  AUDIT-02 &
  L'audit log è accessibile in sola lettura esclusivamente dall'Amministratore.
  Non deve essere possibile modificare o eliminare voci del log dall'interfaccia applicativa. \\[6pt]

  AUDIT-03 &
  Il log deve essere ricercabile per intervallo di date, utente, tipo di operazione. \\

  \bottomrule
\end{longtable}

% ─────────────────────────────────────────────────────────────────────────────
\subsection{Requisiti Non Funzionali}
% ─────────────────────────────────────────────────────────────────────────────

\subsubsection{Prestazioni}

\begin{longtable}{@{}p{2.5cm}p{9.5cm}@{}}
  \toprule
  \textbf{ID} & \textbf{Descrizione} \\
  \midrule
  \endhead

  PERF-01 &
  Il tempo di risposta per le operazioni interattive (caricamento pagine, ricerche, salvataggio
  record) deve essere inferiore a 7 secondi nel 90\% dei casi, con un numero di utenti
  concorrenti fino a 10. \\[6pt]

  PERF-02 &
  Il processo di importazione di un file Excel con 20 righe deve completarsi (validazione
  inclusa) in meno di 30 secondi. \\[6pt]

  PERF-03 &
  Il calcolo mensile delle imposte sull'intera base dati deve completarsi in meno di
  \textbf{[DA DEFINIRE]} minuti (dipende dalla dimensione storica dei dati; da concordare
  dopo analisi dell'algoritmo legacy). \\

  \bottomrule
\end{longtable}

\subsubsection{Affidabilità e disponibilità}

\begin{longtable}{@{}p{2.5cm}p{9.5cm}@{}}
  \toprule
  \textbf{ID} & \textbf{Descrizione} \\
  \midrule
  \endhead

  AVAIL-01 &
  Il sistema deve essere disponibile durante l'orario lavorativo (lun--ven, 08:00--18:00)
  con un uptime target del 99\% (escluse finestre di manutenzione pianificate). \\[6pt]

  AVAIL-02 &
  Il backup del database deve essere eseguito giornalmente (fuori orario lavorativo).
  Il Recovery Point Objective (RPO) è di 24 ore; il Recovery Time Objective (RTO) è
  \textbf{[DA DEFINIRE]}. \\

  \bottomrule
\end{longtable}

\subsubsection{Sicurezza}

\begin{longtable}{@{}p{2.5cm}p{9.5cm}@{}}
  \toprule
  \textbf{ID} & \textbf{Descrizione} \\
  \midrule
  \endhead

  SEC-01 &
  Tutta la comunicazione tra client e server deve avvenire tramite HTTPS (TLS 1.2 o
  superiore). \\[6pt]

  SEC-02 &
  Le password degli utenti devono essere salvate con hashing e salting (es.\ bcrypt o
  Argon2). Non devono mai essere memorizzate in chiaro. \\[6pt]

  SEC-03 &
  I dati personali dei contribuenti e i dati fiscali devono essere trattati in conformità
  al Regolamento UE 2016/679 (GDPR). \\[6pt]

  SEC-04 &
  \textbf{[DA APPROFONDIRE]} Crittografia dei dati a riposo nel database (valutare in
  base alla scelta infrastrutturale cloud vs.\ on-premise). \\

  \bottomrule
\end{longtable}

\subsubsection{Manutenibilità e portabilità}

\begin{longtable}{@{}p{2.5cm}p{9.5cm}@{}}
  \toprule
  \textbf{ID} & \textbf{Descrizione} \\
  \midrule
  \endhead

  MAINT-01 &
  Il codice sorgente deve essere versionato con git e ospitato nel repository del progetto.
  Ogni funzionalità deve essere sviluppata su branch separati con pull request e code review. \\[6pt]

  PORT-01 &
  Il sistema deve poter essere eseguito sia su infrastruttura Azure che su un server
  on-premise Windows/Linux, senza modifiche al codice applicativo (solo configurazione). \\

  \bottomrule
\end{longtable}

\subsubsection{Vincoli generali e conformità}

\begin{longtable}{@{}p{2.5cm}p{9.5cm}@{}}
  \toprule
  \textbf{ID} & \textbf{Descrizione} \\
  \midrule
  \endhead

  CONST-01 &
  L'interfaccia utente deve essere fruibile dai browser: Google Chrome (ultima versione),
  Microsoft Edge (ultima versione), Mozilla Firefox (ultima versione). \\[6pt]

  CONST-02 &
  Il sistema deve prevedere un ambiente di \textit{staging} separato dalla produzione per
  i test pre-rilascio. \\[6pt]

  CONST-03 &
  Il contratto include 6--12 mesi di supporto post-lancio: correzione di bug, assistenza
  agli utenti, aggiornamenti di sicurezza minori. \\

  \bottomrule
\end{longtable}

% =============================================================================
\appendix
\section{Flusso Operativo As-Is e To-Be}
% =============================================================================

\subsection{Flusso As-Is (sistema attuale)}

\begin{enumerate}
  \item Il Catasto invia i file Excel mensilmente.
  \item Un operatore inserisce manualmente i dati nel sistema VB su PC Windows~7.
  \item Il sistema calcola le imposte.
  \item Un responsabile rivede i risultati.
  \item Un operatore genera i PDF e invia le email manualmente da un PC con internet.
\end{enumerate}

\subsection{Flusso To-Be (sistema proposto)}

\begin{enumerate}
  \item Il Catasto invia i file Excel mensilmente.
  \item L'Operatore carica il file nel sistema web (da qualsiasi PC, anche in smart working).
  \item Il sistema valida i dati e li mette in stato «in attesa».
  \item Il Responsabile di Area approva i dati importati.
  \item Il sistema esegue il calcolo delle imposte (schedulato o manuale).
  \item Il sistema genera i PDF di notifica.
  \item Il Responsabile di Area approva i PDF.
  \item Il sistema invia automaticamente le email ai contribuenti.
  \item Tutte le operazioni vengono registrate nell'audit log.
\end{enumerate}

% =============================================================================
\section{Questioni Aperte}
% =============================================================================

Le seguenti decisioni sono ancora da concordare con il cliente o il gruppo di sviluppo:

\begin{enumerate}
  \item \textbf{Scelta infrastrutturale:} Azure cloud vs.\ server on-premise. Impatta su
        sicurezza, backup, autenticazione e portabilità.
  \item \textbf{Integrazione algoritmo legacy:} Chiamata diretta alla DLL VB, riscrittura
        in linguaggio moderno, o esposizione come microservizio. Richiede l'accesso al
        sorgente VB.
  \item \textbf{Strategia di migrazione dati:} Come trasferire lo storico dal sistema legacy
        al nuovo database (one-shot vs.\ migrazione graduale).
  \item \textbf{Formato e schema Excel:} Lo schema del file fornito dal Catasto deve essere
        documentato formalmente per implementare il parser.
  \item \textbf{Dimensione massima allegati:} Limite per i file PDF/ZIP allegati ai record.
  \item \textbf{RTO (Recovery Time Objective):} Tempo massimo di ripristino accettabile
        in caso di disastro.
  \item \textbf{Tempo massimo calcolo mensile:} Da misurare sull'algoritmo legacy con dati
        reali per definire il SLA.
  \item \textbf{Portale pubblico SPID:} Confermarne l'esclusione dalla v1.0 e pianificare
        la v1.1.
\end{enumerate}

\end{document}
